\documentclass[aspectratio=169,utf8,t]{beamer}

% Idioma e Codificacao
\usepackage[brazilian]{babel}
\usepackage{booktabs}
\usepackage{array}
\usepackage[normalem]{ulem}
\usepackage{tikz}
\usetikzlibrary{shapes,shapes.geometric,arrows,arrows.meta,positioning,calc,fit}

% Carrega o Tema Modelo dos Scouts (Widescreen 16:9, Nunito Sans)
\usetheme{scouts}

\definecolor{amber}{RGB}{245, 158, 11}

% Metadados Institucionais
\author{Prof. Me. Reginaldo Fernandes}
\institute{Instituto Federal do Ceará -- Campus Tauá \\ Tecnologia em Telemática}
\title{Banco de Dados}
\subtitle{Aula 08: Laboratório de Modelagem Conceitual com Ferramentas CASE (brModelo e Draw.io)}
\date{20 de Outubro de 2026}

\begin{document}

% ------------------------------------------------------------------------------
% Capa da Apresentacao
% ------------------------------------------------------------------------------
\changecolours[inverse]{ScoutGreen}
\frame[plain]{\titlepage}

% ------------------------------------------------------------------------------
% Sumario Executivo da Aula
% ------------------------------------------------------------------------------
\changecolours{ScoutGreen}
\begin{frame}{Agenda do Encontro}
  \begin{columns}[t]
    \begin{column}{0.48\textwidth}
      \scriptsize
      \tableofcontents[sections={1-5}]
    \end{column}
    \begin{column}{0.48\textwidth}
      \scriptsize
      \tableofcontents[sections={6-9}]
    \end{column}
  \end{columns}
\end{frame}

% ------------------------------------------------------------------------------
% Modulos Tematicos da Aula
% ------------------------------------------------------------------------------
\changecolours{ScoutGreen}
\section[Revisão e Objetivos]{Ponte Pedagógica e Objetivos de Aprendizagem}

% ------------------------------------------------------------------------------
\twocolframe[title={Do Rascunho em Papel (Aula 07) ao Artefato CASE (Aula 08)},leftcol=7.0cm,rightcol=7.0cm]{%
  \alert{O Que Consolidamos na Aula 07 (Sábado)?}\par
  \itemseps[topsepi=0.5mm,itemsepi=1mm,labelsep=2mm,leftmargini=3mm]
  \begin{itemize}
    \item \textbf{Escopo da AVP1:} Narrativa completa das regras de negócio do projeto semestral.
    \item \textbf{Catálogo Completo:} Entidades regulares, atributos de Chen, relacionamentos ($1:1, 1:N, N:M$, recursivo), entidade fraca e hierarquia EER.
    \item \textbf{Esboço Manual:} Desenho preliminar à mão na folha do caderno de estudo dirigido.
  \end{itemize}
}{%
  \alert{A Missão Prática da Aula 08}\par
  \itemseps[topsepi=0.5mm,itemsepi=1mm,labelsep=2mm,leftmargini=3mm]
  \begin{itemize}
    \item \textbf{Digitalização Profissional:} Transformar o esboço manual em um Diagrama Entidade-Relacionamento Estendido (DER/EER) impecável usando \textbf{Ferramentas CASE}.
    \item \textbf{Bilinguismo Visual:} Dominar tanto o padrão internacional de \textbf{Peter Chen (Draw.io)} quanto o padrão brasileiro de \textbf{Heuser (brModelo)}.
    \item \textbf{Dicionário de Dados:} Documentar em tabelas descritivas as regras semânticas do modelo.
  \end{itemize}
}

% ------------------------------------------------------------------------------
\twocolframe[title={Objetivos de Aprendizagem da Aula 08},leftcol=7.0cm,rightcol=7.0cm]{%
  \alert{Competências Conceituais e Notacionais}\par
  \itemseps[topsepi=0.5mm,itemsepi=1mm,labelsep=2mm,leftmargini=3mm]
  \begin{itemize}
    \item \textbf{Ferramentas CASE:} Diferenciar editores gráficos puros de ferramentas baseadas em metamodelo de dados.
    \item \textbf{Tradução Chen $\leftrightarrow$ Heuser:} Converter elipses e linhas duplas para pinos identificadores e pares $(min, max)$.
    \item \textbf{Especialização no brModelo:} Aplicar a dupla de restrições $(t/p, e/c)$ equivalente a $(\text{Total/Parcial}, d/o)$.
  \end{itemize}
}{%
  \alert{Competências Práticas de Modelagem}\par
  \itemseps[topsepi=0.5mm,itemsepi=1mm,labelsep=2mm,leftmargini=3mm]
  \begin{itemize}
    \item \textbf{Operação no brModelo e Draw.io:} Construir diagramas limpos, alinhados em grade e sem cruzamento de arestas.
    \item \textbf{Prevenção de Antipadrões:} Evitar chaves estrangeiras no nível conceitual e entidades desconectadas.
    \item \textbf{Modelagem Descritiva e Textual:} Elaborar o Dicionário de Dados Conceitual e a Matriz de Rastreabilidade.
  \end{itemize}
}

\changecolours{ScoutGreen}
\section[Ferramentas CASE]{O Conceito de Ferramentas CASE em Banco de Dados}

% ------------------------------------------------------------------------------
\twocolframe[title={Engenharia de Software Assistida por Computador (CASE)},leftcol=7.0cm,rightcol=7.0cm]{%
  \alert{O Que é uma Ferramenta CASE?}\par
  \itemseps[topsepi=0.5mm,itemsepi=1mm,labelsep=2mm,leftmargini=3mm]
  \begin{itemize}
    \item \textbf{Sigla Clássica:} \textit{Computer-Aided Software Engineering}.
    \item \textbf{Além do Desenho:} Enquanto um editor comum enxerga apenas formas soltas, uma ferramenta CASE mantém um \textbf{Metamodelo Semântico} subjacente.
    \item \textbf{Funções Essenciais:} Padronização visual (Chen e Heuser), repositório de metadados (Dicionário de Dados) e apoio à transição Conceitual $\to$ Lógico.
  \end{itemize}
}{%
  \alert{Ferramentas Adotadas no Laboratório}\par
  \itemseps[topsepi=0.5mm,itemsepi=1mm,labelsep=2mm,leftmargini=3mm]
  \begin{itemize}
    \item \textbf{1. brModelo (Web / Desktop 3.x):} Criado por Carlos H. Candido seguindo o livro de \textbf{Carlos Alberto Heuser} (UFRGS), com motor CASE nativo para esquemas conceituais e lógicos.
    \item \textbf{2. Draw.io (diagrams.net):} Ferramenta vetorial aberta com a biblioteca \textit{Entity Relation}, ideal para desenhar a \textbf{notação original de Peter Chen (1976)} e \textbf{Elmasri \& Navathe}.
  \end{itemize}
}

% ------------------------------------------------------------------------------
\begin{frame}{Fluxo de Modelagem de Dados com Ferramentas CASE}
  \vspace{0.15cm}
  \begin{center}
    \tikz[scale=0.68, transform shape]{
      \node[draw=ScoutNavy, fill=ScoutNavy!10, rounded corners=4pt, thick, minimum width=3.1cm, minimum height=1.4cm, align=center, font=\sffamily\scriptsize] (req) at (0,0) {
        \textbf{1. Minimundo (Aula 07)}\\
        Narrativa Descritiva\\e Regras de Negócio
      };
      \node[draw=ScoutGreen, fill=ScoutGreen!15, rounded corners=4pt, very thick, minimum width=3.4cm, minimum height=1.4cm, align=center, font=\sffamily\scriptsize] (case) at (4.5,0) {
        \textbf{2. Conceitual (Aula 08)}\\
        DER / EER no \textbf{brModelo}\\ou \textbf{Draw.io} + Dicionário
      };
      \node[draw=ScoutNavy, fill=ScoutNavy!10, rounded corners=4pt, thick, minimum width=3.2cm, minimum height=1.4cm, align=center, font=\sffamily\scriptsize] (log) at (9.0,0) {
        \textbf{3. Lógico (Aulas 09-11)}\\
        Esquema Relacional\\Formal e Chaves (PK/FK)
      };
      \node[draw=amber!90!black, fill=amber!15, rounded corners=4pt, thick, minimum width=3.0cm, minimum height=1.4cm, align=center, font=\sffamily\scriptsize] (norm) at (13.3,0) {
        \textbf{4. Refinamento (Aulas 16-18)}\\
        Dependências Funcionais\\e Normalização (3FN)
      };

      \draw[-{Latex[length=2.3mm]}, very thick, ScoutNavy] (req) -- (case);
      \draw[-{Latex[length=2.3mm]}, very thick, ScoutGreen] (case) -- (log);
      \draw[-{Latex[length=2.3mm]}, very thick, ScoutNavy] (log) -- (norm);
    }
  \end{center}
  \vspace{0.2cm}
  \begin{columns}[t]
    \begin{column}{0.48\textwidth}
      \footnotesize
      \alert{Por Que Modelar Antes de Implementar?}\par
      A modelagem conceitual isola as regras de negócio da complexidade técnica, evitando perda de restrições de participação e ambiguidades estruturais.
    \end{column}
    \begin{column}{0.48\textwidth}
      \footnotesize
      \alert{Qual Ferramenta Escolher na AVP1?}\par
      O aluno pode entregar o DER da AVP1 tanto no \textbf{Draw.io} (notação Peter Chen) quanto no \textbf{brModelo} (notação Heuser), mantendo consistência visual!
    \end{column}
  \end{columns}
\end{frame}

\changecolours{ScoutGreen}
\section[Chen vs. Heuser]{O Duelo de Notações: Peter Chen (Draw.io) vs. Heuser (brModelo)}

% ------------------------------------------------------------------------------
\begin{frame}{Atributos: Peter Chen (Draw.io) vs. Heuser (brModelo)}
  \vspace{-0.15cm}
  \begin{columns}[t]
    \begin{column}{0.49\textwidth}
      \footnotesize\alert{1. Peter Chen Clássico (Draw.io)}\par
      \vspace{0.12cm}
      \begin{center}
        \tikz[scale=0.64, transform shape]{
          \node[draw=ScoutGreen, fill=ScoutGreen!15, thick, minimum width=2.3cm, minimum height=0.7cm, font=\bfseries\sffamily\scriptsize] (ent) at (0,0) {COLABORADOR};
          \node[draw=ScoutNavy, fill=ScoutNavy!10, ellipse, font=\tiny\sffamily\bfseries] (pk) at (-2.5, 1.0) {\underline{matricula}};
          \node[draw=black!80, fill=white, ellipse, font=\tiny\sffamily] (nome) at (-2.5, -0.1) {nome};
          \node[draw=black!80, fill=purple!10, dashed, ellipse, font=\tiny\sffamily] (tempo) at (-2.3, -1.1) {tempo\_serv};
          \node[draw=black!80, fill=pink!15, double, ellipse, font=\tiny\sffamily] (tel) at (2.2, 1.0) {telefones};
          \node[draw=black!80, fill=white, ellipse, font=\tiny\sffamily] (end) at (2.2, -0.4) {endereco};
          \node[draw=black!80, fill=white, ellipse, font=\tiny\sffamily] (rua) at (3.8, 0.1) {rua};
          \node[draw=black!80, fill=white, ellipse, font=\tiny\sffamily] (num) at (3.8, -0.9) {num};

          \draw[semithick] (ent) -- (pk);
          \draw[semithick] (ent) -- (nome);
          \draw[semithick, dashed] (ent) -- (tempo);
          \draw[semithick] (ent) -- (tel);
          \draw[semithick] (ent) -- (end);
          \draw[semithick] (end) -- (rua);
          \draw[semithick] (end) -- (num);
        }
      \end{center}
      \vspace{0.05cm}
      \scriptsize
      $\bullet$ \textbf{Identificador (PK):} Elipse com \underline{sublinhado contínuo}.\\
      $\bullet$ \textbf{Simples:} Elipse de borda contínua simples.\\
      $\bullet$ \textbf{Composto:} Elipse subdividida em elipses menores.\\
      $\bullet$ \textbf{Multivalorado:} Elipse de \textbf{borda dupla} (\texttt{telefones}).\\
      $\bullet$ \textbf{Derivado:} Elipse de \textbf{borda tracejada} (\texttt{tempo\_serv}).
    \end{column}
    \begin{column}{0.49\textwidth}
      \footnotesize\alert{2. Carlos A. Heuser (brModelo)}\par
      \vspace{0.12cm}
      \begin{center}
        \tikz[scale=0.66, transform shape]{
          \node[draw=ScoutGreen, fill=ScoutGreen!15, thick, minimum width=2.4cm, minimum height=2.5cm, font=\bfseries\sffamily\scriptsize] (entb) at (0,0) {COLABORADOR};

          % Pinos estilo brModelo
          \draw[semithick] (1.2, 0.85) -- (1.8, 0.85);
          \filldraw[black] (1.8, 0.85) circle (2.3pt) node[right=3pt, font=\tiny\sffamily\bfseries] {matricula};

          \draw[semithick] (1.2, 0.30) -- (1.8, 0.30);
          \filldraw[draw=black, fill=white, semithick] (1.8, 0.30) circle (2.3pt) node[right=3pt, font=\tiny\sffamily] {nome};

          \draw[semithick] (1.2, -0.25) -- (1.8, -0.25);
          \filldraw[draw=black, fill=white, semithick] (1.8, -0.25) circle (2.3pt) node[right=3pt, font=\tiny\sffamily] {telefones \textbf{(1, n)}};

          \draw[semithick] (1.2, -0.80) -- (1.8, -0.80);
          \filldraw[draw=black, fill=white, semithick] (1.8, -0.80) circle (2.3pt) node[right=3pt, font=\tiny\sffamily] (lblend) {endereco};

          \draw[semithick] (lblend.east) -- ++(0.45, 0.25);
          \filldraw[draw=black, fill=white, semithick] ($(lblend.east)+(0.45,0.25)$) circle (1.9pt) node[right=2pt, font=\tiny\sffamily] {rua};

          \draw[semithick] (lblend.east) -- ++(0.45, -0.25);
          \filldraw[draw=black, fill=white, semithick] ($(lblend.east)+(0.45,-0.25)$) circle (1.9pt) node[right=2pt, font=\tiny\sffamily] {num};
        }
      \end{center}
      \vspace{0.05cm}
      \scriptsize
      $\bullet$ \textbf{Identificador (PK):} Pino com \textbf{círculo preto preenchido} ($\bullet$).\\
      $\bullet$ \textbf{Simples:} Pino com \textbf{círculo branco vazio} ($\circ$).\\
      $\bullet$ \textbf{Composto:} Pino branco ramificado em sub-pinos.\\
      $\bullet$ \textbf{Multivalorado:} Pino branco + cardinalidade \textbf{(0, n)} ou \textbf{(1, n)}.\\
      $\bullet$ \textbf{Derivado:} Especificado no Dicionário de Dados textual.
    \end{column}
  \end{columns}
\end{frame}

% ------------------------------------------------------------------------------
\begin{frame}{Cardinalidades e Participação: Chen vs. Heuser}
  \vspace{-0.15cm}
  \begin{columns}[t]
    \begin{column}{0.47\textwidth}
      \footnotesize\alert{1. Anatomia das Linhas e Pares $(min, max)$}\par
      \vspace{0.15cm}
      \begin{center}
        \tikz[scale=0.72, transform shape]{
          % Chen
          \node[font=\scriptsize\bfseries\sffamily, ScoutNavy] at (0, 1.35) {Padrão Peter Chen / Elmasri (Draw.io):};
          \node[draw=ScoutGreen, fill=ScoutGreen!15, thick, minimum width=1.9cm, minimum height=0.6cm, font=\bfseries\sffamily\tiny] (m1) at (-2.8, 0.45) {MEDICO};
          \node[diamond, aspect=1.8, draw=ScoutNavy, fill=ScoutNavy!12, thick, inner sep=1.2pt, font=\bfseries\sffamily\tiny] (r1) at (0, 0.45) {DIRIGE};
          \node[draw=ScoutGreen, fill=ScoutGreen!15, thick, minimum width=1.9cm, minimum height=0.6cm, font=\bfseries\sffamily\tiny] (d1) at (2.8, 0.45) {DEPTO};
          \draw[semithick] (m1) -- node[above, font=\scriptsize\bfseries, ScoutNavy] {1} (r1);
          \draw[semithick] ($(r1.east)+(0,0.05)$) -- node[above, font=\scriptsize\bfseries, ScoutNavy] {1} ($(d1.west)+(0,0.05)$);
          \draw[semithick] ($(r1.east)+(0,-0.05)$) -- ($(d1.west)+(0,-0.05)$);

          % Callouts Chen
          \node[font=\tiny\sffamily, text=black!85, align=center] (c_simp) at (-1.4, -0.5) {\textbf{Linha Simples}\\Participação \textbf{Parcial}\\(Opcional $\equiv$ mín \textbf{0})};
          \draw[-Stealth, thin, ScoutGreen] (c_simp.north) -- (-1.4, 0.40);

          \node[font=\tiny\sffamily, text=black!85, align=center] (c_dup) at (1.4, -0.5) {\textbf{Linha Dupla}\\Participação \textbf{Total}\\(Obrigatória $\equiv$ mín \textbf{1})};
          \draw[-Stealth, thin, ScoutGreen] (c_dup.north) -- (1.4, 0.38);

          % Heuser / brModelo
          \node[font=\scriptsize\bfseries\sffamily, ScoutNavy] at (0, -1.5) {Padrão Carlos A. Heuser (brModelo):};
          \node[draw=ScoutGreen, fill=ScoutGreen!15, thick, minimum width=1.9cm, minimum height=0.6cm, font=\bfseries\sffamily\tiny] (m2) at (-2.8, -2.35) {MEDICO};
          \node[diamond, aspect=1.8, draw=ScoutNavy, fill=ScoutNavy!12, thick, inner sep=1.2pt, font=\bfseries\sffamily\tiny] (r2) at (0, -2.35) {DIRIGE};
          \node[draw=ScoutGreen, fill=ScoutGreen!15, thick, minimum width=1.9cm, minimum height=0.6cm, font=\bfseries\sffamily\tiny] (d2) at (2.8, -2.35) {DEPTO};
          \draw[semithick] (m2) -- node[above, font=\scriptsize\bfseries, ScoutNavy] {(0, 1)} (r2);
          \draw[semithick] (r2) -- node[above, font=\scriptsize\bfseries, ScoutNavy] {(1, 1)} (d2);

          % Callouts Heuser
          \node[font=\tiny\sffamily, text=black!85, align=center] (h_p1) at (-1.4, -3.25) {\textbf{mín = 0} (Opcional)\\ \textbf{máx = 1} (No máx. 1)};
          \draw[-Stealth, thin, ScoutNavy] (h_p1.north) -- (-1.4, -2.40);

          \node[font=\tiny\sffamily, text=black!85, align=center] (h_p2) at (1.4, -3.25) {\textbf{mín = 1} (Obrigatório)\\ \textbf{máx = 1} (No máx. 1)};
          \draw[-Stealth, thin, ScoutNavy] (h_p2.north) -- (1.4, -2.40);
        }
      \end{center}
    \end{column}
    \begin{column}{0.51\textwidth}
      \footnotesize\alert{2. Guia Completo dos 4 Pares $(min, max)$}\par
      \vspace{0.12cm}
      \renewcommand{\arraystretch}{1.28}
      \resizebox{\linewidth}{!}{%
        \begin{tabular}{@{}ccp{4.2cm}@{}}
          \toprule
          \textbf{brModelo} & \textbf{Draw.io (Chen)} & \textbf{Significado na Regra de Negócio} \\ \midrule
          \textbf{\texttt{(0, 1)}} &
          \tikz[baseline=-2pt]{
            \draw[semithick, ScoutNavy] (0,0) -- node[above=0pt, font=\scriptsize\bfseries]{1} (1.4,0);
          } &
          \textbf{Parcial e Único:} Participa de \textbf{zero} (opcional) ou no máximo \textbf{1} ocorrência. \\
          \textbf{\texttt{(1, 1)}} &
          \tikz[baseline=-2pt]{
            \draw[semithick, ScoutNavy] (0,0.04) -- node[above=0pt, font=\scriptsize\bfseries]{1} (1.4,0.04);
            \draw[semithick, ScoutNavy] (0,-0.04) -- (1.4,-0.04);
          } &
          \textbf{Total e Único:} Participação obrigatória em \textbf{exatamente 1} ocorrência. \\
          \textbf{\texttt{(0, n)}} &
          \tikz[baseline=-2pt]{
            \draw[semithick, ScoutNavy] (0,0) -- node[above=0pt, font=\scriptsize\bfseries]{N} (1.4,0);
          } &
          \textbf{Parcial e Múltiplo:} Participa de \textbf{zero} (opcional), 1 ou \textbf{várias ($N$)} ocorrências. \\
          \textbf{\texttt{(1, n)}} &
          \tikz[baseline=-2pt]{
            \draw[semithick, ScoutNavy] (0,0.04) -- node[above=0pt, font=\scriptsize\bfseries]{N} (1.4,0.04);
            \draw[semithick, ScoutNavy] (0,-0.04) -- (1.4,-0.04);
          } &
          \textbf{Total e Múltiplo:} Participação obrigatória em \textbf{pelo menos 1} ou \textbf{várias ($N$)}. \\ \bottomrule
        \end{tabular}%
      }
      \vspace{0.15cm}
      \scriptsize
      $\bullet$ \textbf{Síntese:} Em Heuser, o 1º valor (\texttt{0} ou \texttt{1}) substitui a linha simples/dupla de Chen, e o 2º valor (\texttt{1} ou \texttt{n}) indica a razão máxima!
    \end{column}
  \end{columns}
\end{frame}

% ------------------------------------------------------------------------------
\begin{frame}{Especialização EER: Círculo $(d/o)$ vs. Triângulo $(t/p, e/c)$}
  \vspace{-0.15cm}
  \begin{columns}[t]
    \begin{column}{0.48\textwidth}
      \footnotesize\alert{1. Anatomia dos Símbolos e Siglas}\par
      \vspace{0.08cm}
      \begin{center}
        \tikz[scale=0.62, transform shape]{
          % EER Chen
          \node[font=\scriptsize\bfseries\sffamily, ScoutNavy] at (-2.1, 1.50) {Draw.io (Elmasri / Chen)};
          \node[draw=ScoutGreen, fill=ScoutGreen!15, thick, minimum width=2.1cm, minimum height=0.52cm, font=\bfseries\sffamily\tiny] (s1) at (-2.1, 0.80) {COLABORADOR};
          \node[draw=ScoutNavy, fill=amber!25, circle, thick, inner sep=1.8pt, font=\bfseries\sffamily\scriptsize] (c1) at (-2.1, -0.25) {d};
          \node[draw=ScoutGreen, fill=ScoutGreen!15, thick, minimum width=1.4cm, minimum height=0.48cm, font=\bfseries\sffamily\tiny] (b1) at (-3.05, -1.25) {MEDICO};
          \node[draw=ScoutGreen, fill=ScoutGreen!15, thick, minimum width=1.4cm, minimum height=0.48cm, font=\bfseries\sffamily\tiny] (b2) at (-1.15, -1.25) {ATEND};
          \draw[semithick] ($(s1.south)+(-0.04,0)$) -- ($(c1.north)+(-0.04,0)$);
          \draw[semithick] ($(s1.south)+(0.04,0)$) -- ($(c1.north)+(0.04,0)$);
          \draw[semithick] (c1) -- (b1);
          \draw[semithick] (c1) -- (b2);

          % Callouts Chen EER
          \node[font=\tiny\sffamily, align=right, text=black!85] (cl_dup) at (-4.15, 0.32) {\textbf{Linha Dupla} = Total\\(\textbf{Simples} = Parcial)};
          \draw[-Stealth, thin, ScoutGreen] (cl_dup.east) -- (-2.2, 0.32);

          \node[font=\tiny\sffamily, align=right, text=black!85] (cl_cir) at (-4.15, -0.35) {\textbf{Letra no Círculo:}\\\textbf{d} = Disjunta | \textbf{o} = Sobreposta};
          \draw[-Stealth, thin, ScoutNavy] (cl_cir.east) -- (c1.west);

          % EER brModelo
          \node[font=\scriptsize\bfseries\sffamily, ScoutNavy] at (1.9, 1.50) {brModelo (Heuser)};
          \node[draw=ScoutGreen, fill=ScoutGreen!15, thick, minimum width=2.1cm, minimum height=0.52cm, font=\bfseries\sffamily\tiny] (s2) at (1.9, 0.80) {COLABORADOR};
          \node[regular polygon, regular polygon sides=3, draw=ScoutNavy, fill=amber!25, thick, inner sep=2.4pt] (t2) at (1.9, -0.25) {};
          \node[font=\scriptsize\bfseries\sffamily, ScoutNavy, right=2pt of t2] (lblte) {(t, e)};
          \node[draw=ScoutGreen, fill=ScoutGreen!15, thick, minimum width=1.4cm, minimum height=0.48cm, font=\bfseries\sffamily\tiny] (h1) at (0.95, -1.25) {MEDICO};
          \node[draw=ScoutGreen, fill=ScoutGreen!15, thick, minimum width=1.4cm, minimum height=0.48cm, font=\bfseries\sffamily\tiny] (h2) at (2.85, -1.25) {ATEND};
          \draw[semithick] (s2) -- (t2);
          \draw[semithick] (t2) -- (h1);
          \draw[semithick] (t2) -- (h2);
        }
      \end{center}
      \vspace{-0.12cm}
      \scriptsize
      $\bullet$ \textbf{1ª Letra no brModelo (\texttt{t} ou \texttt{p} — Completude):}\\
      \hspace{0.15cm}-- \textbf{\texttt{t} (Total):} Todo registro pertence a alguma subclasse.\\
      \hspace{0.15cm}-- \textbf{\texttt{p} (Parcial):} Pode não pertencer a nenhuma subclasse.\\
      $\bullet$ \textbf{2ª Letra no brModelo (\texttt{e} ou \texttt{c} — Disjunção):}\\
      \hspace{0.15cm}-- \textbf{\texttt{e} (Exclusiva $\equiv$ \texttt{d} Disjunta):} Pertence a \textbf{uma única} subclasse.\\
      \hspace{0.15cm}-- \textbf{\texttt{c} (Compartilhada $\equiv$ \texttt{o} Sobreposta):} Pode pertencer a \textbf{ambas}.
    \end{column}
    \begin{column}{0.50\textwidth}
      \footnotesize\alert{2. As 4 Combinações de Herança EER}\par
      \vspace{0.12cm}
      \renewcommand{\arraystretch}{1.28}
      \resizebox{\linewidth}{!}{%
        \begin{tabular}{@{}ccp{4.1cm}@{}}
          \toprule
          \textbf{Draw.io (Chen)} & \textbf{brModelo (Heuser)} & \textbf{Significado Completo da Sigla} \\ \midrule
          \tikz[baseline=-4pt]{
            \draw[semithick] (-0.04,0.35) -- (-0.04,0.08);
            \draw[semithick] (0.04,0.35) -- (0.04,0.08);
            \node[draw=ScoutNavy, fill=amber!25, circle, inner sep=1.4pt, font=\scriptsize\bfseries] at (0,-0.12) {d};
          } &
          \tikz[baseline=-3pt]{
            \node[regular polygon, regular polygon sides=3, draw=ScoutNavy, fill=amber!25, inner sep=2.0pt] (tr) at (0,0) {};
            \node[font=\scriptsize\bfseries, ScoutNavy, right=2pt of tr] {(t, e)};
          } &
          \textbf{Total e Exclusiva (Disjunta):} Todo colaborador \textbf{deve} estar em uma subclasse, e em \textbf{apenas uma} (\texttt{d}). \\
          \tikz[baseline=-4pt]{
            \draw[semithick] (0,0.35) -- (0,0.08);
            \node[draw=ScoutNavy, fill=amber!25, circle, inner sep=1.4pt, font=\scriptsize\bfseries] at (0,-0.12) {d};
          } &
          \tikz[baseline=-3pt]{
            \node[regular polygon, regular polygon sides=3, draw=ScoutNavy, fill=amber!25, inner sep=2.0pt] (tr) at (0,0) {};
            \node[font=\scriptsize\bfseries, ScoutNavy, right=2pt of tr] {(p, e)};
          } &
          \textbf{Parcial e Exclusiva (Disjunta):} Opcional; mas se pertencer a subclasse, será a \textbf{apenas uma} (\texttt{d}). \\
          \tikz[baseline=-4pt]{
            \draw[semithick] (-0.04,0.35) -- (-0.04,0.08);
            \draw[semithick] (0.04,0.35) -- (0.04,0.08);
            \node[draw=ScoutNavy, fill=amber!25, circle, inner sep=1.4pt, font=\scriptsize\bfseries] at (0,-0.12) {o};
          } &
          \tikz[baseline=-3pt]{
            \node[regular polygon, regular polygon sides=3, draw=ScoutNavy, fill=amber!25, inner sep=2.0pt] (tr) at (0,0) {};
            \node[font=\scriptsize\bfseries, ScoutNavy, right=2pt of tr] {(t, c)};
          } &
          \textbf{Total e Compartilhada (Sobreposta):} Obrigatório estar em pelo menos uma, podendo \textbf{acumular ambas} (\texttt{o}). \\
          \tikz[baseline=-4pt]{
            \draw[semithick] (0,0.35) -- (0,0.08);
            \node[draw=ScoutNavy, fill=amber!25, circle, inner sep=1.4pt, font=\scriptsize\bfseries] at (0,-0.12) {o};
          } &
          \tikz[baseline=-3pt]{
            \node[regular polygon, regular polygon sides=3, draw=ScoutNavy, fill=amber!25, inner sep=2.0pt] (tr) at (0,0) {};
            \node[font=\scriptsize\bfseries, ScoutNavy, right=2pt of tr] {(p, c)};
          } &
          \textbf{Parcial e Compartilhada (Sobreposta):} Opcional, podendo pertencer a \textbf{múltiplas subclasses} (\texttt{o}). \\ \bottomrule
        \end{tabular}%
      }
    \end{column}
  \end{columns}
\end{frame}

\changecolours{ScoutGreen}
\section[Guia brModelo]{Laboratório Passo a Passo no brModelo (Notação Heuser)}

% ------------------------------------------------------------------------------
\begin{frame}{Operando o brModelo Web / Desktop 3.x: Guia Visual de Formas}
  \vspace{-0.15cm}
  \begin{center}
    \tikz[scale=0.68, transform shape]{
      % Entidade Forte PACIENTE no brModelo
      \node[draw=ScoutGreen, fill=ScoutGreen!15, thick, minimum width=2.2cm, minimum height=1.2cm, font=\bfseries\sffamily\scriptsize] (pac) at (0, 0) {PACIENTE};
      \draw[semithick] (-1.1, 0.35) -- (-1.6, 0.35);
      \filldraw[black] (-1.6, 0.35) circle (2.2pt) node[left=2pt, font=\tiny\sffamily\bfseries] {id\_prontuario};
      \draw[semithick] (-1.1, -0.25) -- (-1.6, -0.25);
      \filldraw[draw=black, fill=white, semithick] (-1.6, -0.25) circle (2.2pt) node[left=2pt, font=\tiny\sffamily] {nome\_pac};

      % Relacionamento Identificador POSSUI_DEP
      \node[diamond, aspect=1.9, draw=ScoutNavy, double, double distance=1.2pt, fill=ScoutNavy!15, thick, inner sep=1.2pt, font=\bfseries\sffamily\tiny] (rel) at (3.4, 0) {POSSUI\_DEP};

      % Entidade Fraca DEPENDENTE no brModelo
      \node[draw=ScoutGreen, double, double distance=1.3pt, fill=ScoutGreen!15, thick, minimum width=2.3cm, minimum height=1.2cm, font=\bfseries\sffamily\scriptsize] (dep) at (6.9, 0) {DEPENDENTE};
      \draw[semithick] (8.05, 0.35) -- (8.55, 0.35);
      \filldraw[draw=black, fill=amber!80, semithick] (8.55, 0.35) circle (2.2pt) node[right=2pt, font=\tiny\sffamily\bfseries] {\dashuline{nome\_dep}};
      \draw[semithick] (8.05, -0.25) -- (8.55, -0.25);
      \filldraw[draw=black, fill=white, semithick] (8.55, -0.25) circle (2.2pt) node[right=2pt, font=\tiny\sffamily] {parentesco};

      \draw[semithick] (pac) -- node[above, font=\tiny\bfseries] {(0, n)} (rel);
      \draw[semithick] (rel) -- node[above, font=\tiny\bfseries] {(1, 1)} (dep);
    }
  \end{center}
  \vspace{0.05cm}
  \begin{columns}[t]
    \begin{column}{0.48\textwidth}
      \scriptsize
      \alert{1. Entidades e Pinos de Atributos:}\par
      $\bullet$ \textbf{Atributo Identificador ($\bullet$):} Clique na entidade e selecione o pino preto preenchido para a chave primária.\\
      $\bullet$ \textbf{Composto e Multivalorado:} Adicione sub-pinos ao atributo composto ou ajuste a cardinalidade para \texttt{(1, n)}.
    \end{column}
    \begin{column}{0.48\textwidth}
      \scriptsize
      \alert{2. Auto-Relacionamento, Fraca e EER:}\par
      $\bullet$ \textbf{Auto-Relacionamento:} Ligue a entidade a si mesma e preencha o campo \textit{Papel} nas duas arestas.\\
      $\bullet$ \textbf{Especialização:} Use o \textbf{Triângulo} e defina no painel lateral \textit{Total/Parcial} (\texttt{t/p}) e \textit{Exclusiva/Compartilhada} (\texttt{e/c}).
    \end{column}
  \end{columns}
\end{frame}

\changecolours{ScoutGreen}
\section[Guia Draw.io]{Laboratório Passo a Passo no Draw.io (Notação Peter Chen)}

% ------------------------------------------------------------------------------
\begin{frame}{Paleta Visual \textit{Entity Relation} de Peter Chen no Draw.io}
  \vspace{-0.15cm}
  \begin{center}
    \tikz[scale=0.66, transform shape]{
      % Linha 1 de formas geometricas reais de Chen
      \node[draw=ScoutGreen, fill=ScoutGreen!15, thick, minimum width=2.3cm, minimum height=0.68cm, font=\bfseries\sffamily\tiny] (f1) at (0, 1.0) {ENTIDADE FORTE};
      \node[font=\tiny\sffamily, below=1pt of f1] {\textit{Entity}};

      \node[draw=ScoutGreen, double, double distance=1.5pt, fill=ScoutGreen!15, thick, minimum width=2.3cm, minimum height=0.68cm, font=\bfseries\sffamily\tiny] (f2) at (3.5, 1.0) {ENTIDADE FRACA};
      \node[font=\tiny\sffamily, below=1pt of f2] {\textit{Weak Entity}};

      \node[diamond, aspect=2.1, draw=ScoutNavy, fill=ScoutNavy!12, thick, inner sep=3.5pt, font=\bfseries\sffamily\tiny] (f3) at (7.0, 1.0) {RELACIONA};
      \node[font=\tiny\sffamily, below=1pt of f3] {\textit{Relationship}};

      \node[diamond, aspect=2.1, draw=ScoutNavy, double, double distance=1.5pt, fill=ScoutNavy!15, thick, inner sep=3.5pt, font=\bfseries\sffamily\tiny] (f4) at (10.6, 1.0) {IDENTIFICA};
      \node[font=\tiny\sffamily, below=1pt of f4] {\textit{Identifying Rel.}};

      % Linha 2 de elipses reais de Chen
      \node[draw=ScoutNavy, fill=ScoutNavy!10, ellipse, thick, minimum width=2.1cm, minimum height=0.58cm, font=\tiny\sffamily\bfseries] (a1) at (0, -0.65) {\underline{atributo\_pk}};
      \node[font=\tiny\sffamily, below=1pt of a1] {\textit{Key Attribute}};

      \node[draw=ScoutNavy, fill=amber!15, ellipse, thick, minimum width=2.1cm, minimum height=0.58cm, font=\tiny\sffamily\bfseries] (a2) at (3.5, -0.65) {\dashuline{discriminador}};
      \node[font=\tiny\sffamily, below=1pt of a2] {\textit{Weak Key}};

      \node[draw=black!80, fill=pink!15, double, double distance=1.4pt, ellipse, thick, minimum width=2.1cm, minimum height=0.58cm, font=\tiny\sffamily] (a3) at (7.0, -0.65) {multivalorado};
      \node[font=\tiny\sffamily, below=1pt of a3] {\textit{Multivalued}};

      \node[draw=black!80, fill=purple!10, dashed, ellipse, thick, minimum width=2.1cm, minimum height=0.58cm, font=\tiny\sffamily] (a4) at (10.6, -0.65) {derivado};
      \node[font=\tiny\sffamily, below=1pt of a4] {\textit{Derived Attribute}};
    }
  \end{center}
  \vspace{0.05cm}
  \begin{columns}[t]
    \begin{column}{0.48\textwidth}
      \scriptsize
      \alert{1. Ativando a Biblioteca no Draw.io:}\par
      $\bullet$ Em \texttt{app.diagrams.net}, clique em \textit{+ Mais Formas} e habilite \textbf{Entity Relation}.\\
      $\bullet$ Arraste diretamente as formas geométricas acima para a página em grade (\textit{Grid}).
    \end{column}
    \begin{column}{0.48\textwidth}
      \scriptsize
      \alert{2. Participação Total e Exportação:}\par
      $\bullet$ \textbf{Linha Dupla:} Selecione a linha entre a entidade e o losango e mude o estilo do traço para \textbf{Double Line}.\\
      $\bullet$ \textbf{Exportação Nitidez Máxima:} Use \textit{Arquivo $\to$ Exportar como $\to$ PDF (Cortar)} ou \textit{PNG (300 DPI)}.
    \end{column}
  \end{columns}
\end{frame}

\changecolours{ScoutGreen}
\section[Boas Práticas e Dicionário]{Padrões de Qualidade Estética, Antipadrões e Dicionário de Dados}

% ------------------------------------------------------------------------------
\begin{frame}{Contraste Visual: Antipadrões Proibidos vs. Modelagem Correta}
  \vspace{-0.15cm}
  \begin{columns}[t]
    \begin{column}{0.49\textwidth}
      \alert{ERRADO: Antipadrões Comuns no DER}\par
      \vspace{0.1cm}
      \begin{center}
        \tikz[scale=0.62, transform shape]{
          % Erro 1: Ligar entidades sem losango e colocar FK como atributo
          \node[draw=red!80!black, fill=red!10, thick, minimum width=1.9cm, minimum height=0.6cm, font=\bfseries\sffamily\tiny] (e1) at (-1.8, 0) {MEDICO};
          \node[draw=red!80!black, fill=red!10, thick, minimum width=1.9cm, minimum height=0.6cm, font=\bfseries\sffamily\tiny] (e2) at (1.8, 0) {DEPTO};
          \draw[very thick, red!80!black] (e1) -- node[above, font=\tiny\bfseries, red!80!black] {Sem Losango!} (e2);

          \node[draw=red!80!black, fill=red!15, ellipse, thick, font=\tiny\sffamily\bfseries] (fk) at (-1.8, 1.1) {cod\_depto\_fk};
          \draw[semithick, red!80!black] (e1) -- (fk);
        }
      \end{center}
      \scriptsize
      $\bullet$ \textbf{1. Atributo de Chave Estrangeira (\texttt{fk}):} Jamais crie atributos como \texttt{cod\_depto\_fk} no DER conceitual!\\
      $\bullet$ \textbf{2. Retângulos Ligados Diretamente:} Duas entidades nunca se ligam sem um losango intermediário.\\
      $\bullet$ \textbf{3. Mistura de Notações:} Nunca misture Pé de Galinha (\textit{Crow's Foot}) com elipses de Chen.
    \end{column}
    \begin{column}{0.49\textwidth}
      \alert{CORRETO: Gramática Canônica Limpa}\par
      \vspace{0.1cm}
      \begin{center}
        \tikz[scale=0.62, transform shape]{
          \node[draw=ScoutGreen, fill=ScoutGreen!15, thick, minimum width=1.9cm, minimum height=0.6cm, font=\bfseries\sffamily\tiny] (c1) at (-2.4, 0) {MEDICO};
          \node[diamond, aspect=1.8, draw=ScoutNavy, fill=ScoutNavy!12, thick, inner sep=1.2pt, font=\bfseries\sffamily\tiny] (rel) at (0, 0) {LOTADO};
          \node[draw=ScoutGreen, fill=ScoutGreen!15, thick, minimum width=1.9cm, minimum height=0.6cm, font=\bfseries\sffamily\tiny] (c2) at (2.4, 0) {DEPTO};

          \node[draw=ScoutNavy, fill=ScoutNavy!10, ellipse, font=\tiny\sffamily\bfseries] (pk1) at (-2.4, 1.1) {\underline{matricula}};
          \node[draw=ScoutNavy, fill=ScoutNavy!10, ellipse, font=\tiny\sffamily\bfseries] (pk2) at (2.4, 1.1) {\underline{cod\_depto}};

          \draw[semithick] (c1) -- (pk1);
          \draw[semithick] (c2) -- (pk2);
          \draw[thick, double, double distance=1.3pt] (c1) -- node[above, font=\tiny\bfseries] {N} (rel);
          \draw[thick, double, double distance=1.3pt] (rel) -- node[above, font=\tiny\bfseries] {1} (c2);
        }
      \end{center}
      \scriptsize
      $\bullet$ \textbf{Nomenclatura Padronizada:} Entidades no singular em \texttt{CAIXA\_ALTA}, verbos ativos nos losangos e atributos em \texttt{snake\_case}.\\
      $\bullet$ \textbf{Topologia Ortogonal:} Entidades alinhadas em grade sem cruzamento de linhas sobre caixas.
    \end{column}
  \end{columns}
\end{frame}

% ------------------------------------------------------------------------------
\twocolframe[title={O Dicionário de Dados Conceitual Complementar},leftcol=7.0cm,rightcol=7.0cm]{%
  \alert{O Que o Desenho Gráfico Não Mostra?}\par
  \itemseps[topsepi=0.5mm,itemsepi=1mm,labelsep=2mm,leftmargini=3mm]
  \begin{itemize}
    \item \textbf{Limitação Visual:} O diagrama DER/EER não detalha a semântica de negócio, o formato de domínio ou a fórmula exata de um atributo derivado.
    \item \textbf{Papel do Dicionário de Dados:} Especificação descritiva e tabular que acompanha obrigatoriamente o DER no projeto da AVP1.
  \end{itemize}
}{%
  \alert{Elementos Obrigatórios do Dicionário}\par
  \itemseps[topsepi=0.5mm,itemsepi=1mm,labelsep=2mm,leftmargini=3mm]
  \begin{itemize}
    \item \textbf{1. Catálogo de Entidades:} Definição clara do papel de cada entidade forte, fraca e subclasse.
    \item \textbf{2. Especificação de Atributos:} Classificação (PK, simples, composto, multivalorado, derivado), obrigatoriedade e regra de cálculo.
    \item \textbf{3. Regras de Negócio Textuais:} Restrições semânticas (ex: ``um médico não supervisiona a si mesmo'').
  \end{itemize}
}

\changecolours{ScoutGreen}
\section[Modelagem Descritiva]{Modelagem Descritiva e Textual: O Dicionário de Dados da AVP1}

% ------------------------------------------------------------------------------
\begin{frame}{Dicionário de Dados: Quadro Descritivo de Atributos}
  \vspace{-0.15cm}
  \footnotesize\alert{Especificação Tabular de Atributos (Estudo de Caso Clínica ``VidaPlena''):}\par
  \vspace{0.12cm}
  \renewcommand{\arraystretch}{1.42}
  \resizebox{\textwidth}{!}{%
    \begin{tabular}{@{}lllp{3.6cm}p{4.8cm}@{}}
      \toprule
      \textbf{Entidade} & \textbf{Natureza} & \textbf{Atributo} & \textbf{Classificação (Chen / Heuser)} & \textbf{Descrição Semântica e Regra de Domínio} \\
      \midrule
      \texttt{COLABORADOR} & Superclasse & \underline{\texttt{matricula}} & Identificador (PK / Pino $\bullet$) & Código funcional inteiro positivo, único e obrigatório. \\
      \texttt{COLABORADOR} & Superclasse & \texttt{endereco} & Composto (\texttt{rua, num, cep}) & Endereço residencial decomposto em subatributos atômicos. \\
      \texttt{COLABORADOR} & Superclasse & \texttt{telefones} & Multivalorado -- \texttt{(1, n)} & Conjunto de números de telefone de contato (mínimo 1). \\
      \texttt{COLABORADOR} & Superclasse & \texttt{tempo\_serv} & Derivado (Elipse Tracejada) & Calculado dinamicamente: $\text{Data Atual} - \text{data\_admissao}$. \\
      \texttt{DEPENDENTE} & Ent. Fraca & \dashuline{\texttt{nome\_dep}} & Discriminador (Chave Parcial) & Identifica o dependente junto ao \underline{\texttt{id\_prontuario}} do paciente titular. \\
      \bottomrule
    \end{tabular}%
  }
  \vspace{0.18cm}
  \scriptsize
  $\bullet$ \textbf{Regra de Documentação da AVP1:} Todo atributo composto, multivalorado, derivado ou chave parcial desenhado na ferramenta CASE deve possuir sua regra de domínio detalhada no Quadro Descritivo acima.
\end{frame}

% ------------------------------------------------------------------------------
\begin{frame}{Dicionário de Dados: Matriz Textual de Relacionamentos}
  \vspace{-0.15cm}
  \footnotesize\alert{Especificação de Associações, Papéis e Cardinalidades (Clínica ``VidaPlena''):}\par
  \vspace{0.12cm}
  \renewcommand{\arraystretch}{1.45}
  \resizebox{\textwidth}{!}{%
    \begin{tabular}{@{}lcp{3.6cm}p{3.8cm}p{3.8cm}@{}}
      \toprule
      \textbf{Relacionamento} & \textbf{Razão} & \textbf{Entidade 1 [Papel] $(min, max)$} & \textbf{Entidade 2 [Papel] $(min, max)$} & \textbf{Atributos Próprios e Observações} \\
      \midrule
      \texttt{DIRIGE} & $1:1$ & \texttt{MEDICO} [diretor] \textbf{\texttt{(0, 1)}} & \texttt{DEPTO\_CLINICO} [dirigido] \textbf{\texttt{(1, 1)}} & \textit{Nenhum} (Participação Total do Departamento). \\
      \texttt{SUPERVISIONA} & $1:N$ & \texttt{MEDICO} [preceptor] \textbf{\texttt{(0, n)}} & \texttt{MEDICO} [residente] \textbf{\texttt{(0, 1)}} & \textit{Nenhum} (Auto-relacionamento recursivo com papéis). \\
      \texttt{POSSUI\_DEP} & $1:N$ & \texttt{PACIENTE} [titular] \textbf{\texttt{(0, n)}} & \texttt{DEPENDENTE} [dep.] \textbf{\texttt{(1, 1)}} & \textit{Nenhum} (Relacionamento Identificador de Entidade Fraca). \\
      \texttt{REALIZA\_EXAME} & $N:M$ & \texttt{PACIENTE} [examinado] \textbf{\texttt{(0, n)}} & \texttt{EXAME} [prescrito] \textbf{\texttt{(0, n)}} & \texttt{data\_solicitacao}, \texttt{justificativa\_clinica}. \\
      \bottomrule
    \end{tabular}%
  }
  \vspace{0.18cm}
  \scriptsize
  $\bullet$ \textbf{Consistência Gráfica $\leftrightarrow$ Textual:} Os pares \texttt{(min, max)} e os papéis documentados nesta matriz devem coincidir exatamente com as linhas e rótulos desenhados no brModelo ou Draw.io.
\end{frame}

\changecolours{ScoutGreen}
\section[Conclusão e Próximos Passos]{Checklist do Entregável CASE e Ponte para a Aula 09}

% ------------------------------------------------------------------------------
\twocolframe[title={Checklist Final de Qualidade do DER/EER (AVP1)},leftcol=7.0cm,rightcol=7.0cm]{%
  \alert{Antes de Exportar o PDF/PNG, Verifique:}\par
  \itemseps[topsepi=0.5mm,itemsepi=1mm,labelsep=2mm,leftmargini=3mm]
  \begin{itemize}
    \item[$\boxtimes$] Todas as entidades fortes possuem chave primária sublinhada ($\underline{\text{pk}}$) ou pino preto ($\bullet$)?
    \item[$\boxtimes$] A entidade fraca possui retângulo duplo, losango duplo, linha dupla e discriminador tracejado?
    \item[$\boxtimes$] A hierarquia EER explicita as duas restrições ($d/o$ e Total/Parcial ou $t/p, e/c$)?
    \item[$\boxtimes$] O auto-relacionamento possui os papéis (\textit{roles}) escritos nas duas arestas?
    \item[$\boxtimes$] Não há nenhuma chave estrangeira desenhada como atributo no DER?
  \end{itemize}
}{%
  \alert{Ponte Pedagógica: Rumo à Aula 09!}\par
  \itemseps[topsepi=0.5mm,itemsepi=1mm,labelsep=2mm,leftmargini=3mm]
  \begin{itemize}
    \item \textbf{Fechamento do Ciclo Conceitual:} Com a Aula 08, concluímos a modelagem conceitual (Unidade III).
    \item \textbf{Aula 09 (21/10/2026) -- O Modelo Relacional (Edgar F. Codd, 1970):} Estudaremos domínios atômicos, tuplas, esquemas de relação, superchaves, chaves candidatas, chave primária (PK), chave estrangeira (FK) e as regras formais de integridade!
  \end{itemize}
}

\changecolours{ScoutGreen}
\section[Referências]{Referências Bibliográficas Recomendadas}

% ------------------------------------------------------------------------------
\begin{frame}{Referências Bibliográficas}
  \itemseps[topsepi=0.5mm,itemsepi=1mm,labelsep=2mm,leftmargini=3.5mm]
  \begin{itemize}
    \scriptsize
    \item \textbf{HEUSER, Carlos Alberto.} \textit{Projeto de Banco de Dados}. 6. ed. Porto Alegre: Bookman, 2009. (Capítulos 2 e 3: Notação Conceitual no padrão brModelo e Dicionário de Dados).
    \item \textbf{CHEN, Peter Pin-Shan.} \textit{The Entity-Relationship Model: Toward a Unified View of Data}. ACM Transactions on Database Systems (TODS), New York, v. 1, n. 1, p. 9--36, Mar. 1976.
    \item \textbf{ELMASRI, Ramez; NAVATHE, Shamkant B.} \textit{Sistemas de Banco de Dados}. 7. ed. São Paulo: Pearson Education do Brasil, 2018. (Capítulos 3 e 4: Modelagem Conceitual e EER).
    \item \textbf{SILBERSCHATZ, Abraham; KORTH, Henry F.; SUDARSHAN, S.} \textit{Sistema de Banco de Dados}. 7. ed. Rio de Janeiro: Elsevier, 2020.
    \item \textbf{CANDIDO, Carlos H.} \textit{brModelo Web \& Desktop: Ferramenta CASE para Modelagem de Bancos de Dados}. Disponível em: \texttt{https://www.brmodeloweb.com}.
  \end{itemize}
  \vspace{0.15cm}
  \begin{center}
    \tiny\color{gray}
    Instituto Federal de Educação, Ciência e Tecnologia do Ceará -- Campus Tauá\\
    Curso Superior de Tecnologia em Telemática
  \end{center}
\end{frame}


\end{document}
